\documentclass[11pt,a4paper]{article}
\usepackage[T1]{fontenc}
\usepackage{lmodern,amsmath,amssymb,amsthm,mathtools,booktabs,microtype}
\usepackage[margin=25mm]{geometry}
\usepackage[colorlinks=true,linkcolor=blue,citecolor=blue,urlcolor=blue]{hyperref}
\setlength{\emergencystretch}{3em}
\newtheorem{theorem}{Theorem}[section]
\newtheorem{lemma}[theorem]{Lemma}
\newtheorem{proposition}[theorem]{Proposition}
\newtheorem{corollary}[theorem]{Corollary}
\theoremstyle{definition}\newtheorem{definition}[theorem]{Definition}
\newcommand{\Q}{\mathbb Q}\newcommand{\Z}{\mathbb Z}
\DeclareMathOperator{\N}{N}\DeclareMathOperator{\Tr}{Tr}
\title{Unit guards and a finite square-complement reduction for elliptic detectors}
\author{}\date{16 November 2026}
\begin{document}\maketitle
\begin{abstract}
For an actual rational specialization of the elliptic point algebra,
every geometric character takes values in $\Q(\sqrt2)$. This implies
that a prepared detector input has nonzero complement $17(17-b)$ at
every character: vanishing would contradict its square class at~17.
Thus rational square complements, if realized by actual detector
tuples, can be encoded with the existing single finite-algebra
coordinate and no additional complement-guard avoidance requirement.
The resulting conditional counts are 21 variables for the general
source and nine for either conditional source, including the input.
The required simultaneous rational-complement realization is stated
explicitly and is not proved. Separately, an elementary construction
gives infinitely many prepared free coefficient inputs at a common
pole parameter; it imposes no native trace realization.
\end{abstract}

\section{The point algebra and its splitting field}
Put $K=\Q(\sqrt2)$. For retained rational coordinates $h_0,h_1$, write
\[
h=h_0+h_1\sqrt2,\qquad h^4-(3-\sqrt2)^2=A+B\sqrt2.
\]
The rank-four point algebra of~\cite{Unit} has basis $1,U,V,W$ and
multiplication table
\begin{equation}\label{eq:block}
\begin{gathered}
U^2=2A+2W,\qquad V^2=A-W,\qquad UV=2B,\\
UW=AU-2BV,\qquad VW=BU-AV,\qquad W^2=A^2-2B^2.
\end{gathered}
\end{equation}
It is etale at every rational pair $(h_0,h_1)$, and its rational
characters correspond to $j=u+v\sqrt2\in K$ satisfying
$j^2=A+B\sqrt2$, with $U=2u$, $V=2v$, and $W=u^2-2v^2$.
An \emph{actual point specialization} means such a rational character
has been chosen for every point block. The corresponding elliptic
point is $(h^2,hj)$.

\begin{lemma}\label{lem:split}
At an actual point specialization, the tensor product $\mathcal B$ of
any finite number of point blocks is split by $K$. In particular every
geometric character of $\mathcal B$ has image in $K$.
\end{lemma}
\begin{proof}
For one block let $j=u+v\sqrt2$ be the chosen root and let $\bar j$ be
its conjugate. Its four geometric characters are
\[
U=\epsilon_+j+\epsilon_-\bar j,\qquad
V=\frac{\epsilon_+j-\epsilon_-\bar j}{\sqrt2},\qquad
W=\epsilon_+\epsilon_-\N_{K/\Q}(j),
\quad \epsilon_+,\epsilon_-\in\{1,-1\}.
\]
Substitution verifies~\eqref{eq:block}. These are four distinct
characters, because $j,\bar j$ are nonzero; they exhaust the
characters of the rank-four etale algebra. Every value lies in $K$.
Tensoring the resulting splitting isomorphisms over $K$ proves the
assertion for $\mathcal B$.
\end{proof}

\begin{lemma}\label{lem:units}
Let $\mathcal B$ be the point algebra at an actual point
specialization. Let $w\in\Q^\times$ have even valuation at~17, and
let $S\in\mathcal B^\times$. If $b=wS^2$, then
\[
                         17(17-b)\in\mathcal B^\times.
\]
\end{lemma}
\begin{proof}
If a geometric character $\chi$ annihilated this element, then
$w\chi(S)^2=17$, and hence
\[
                         \frac w{17}=\chi(S)^{-2}\in K^{\times2}.
\]
A nonzero rational number that is a square in $K$ belongs to
$\Q^{\times2}$ or $2\Q^{\times2}$. Indeed, the irrational component
of $(a+b\sqrt2)^2$ vanishes only when $ab=0$. Both rational square
classes have even valuation at~17, whereas $v_{17}(w/17)$ is odd.
This is a contradiction. A finite etale algebra element is a unit
exactly when every geometric character is nonzero on it.
\end{proof}

\section{The native complement condition}
Both sources in~\cite{Unit} have detector inputs of the form
\begin{equation}\label{eq:detector}
b_i=w_i\left(\frac{\N(\mathcal W_i)}{(\mathcal W_i)_0}\right)^2
    =\frac{4w_i}{\Tr(\mathcal W_i^{-1})^2},\qquad i=0,1.
\end{equation}
For the general source $w_i=\rho R_i$. For a conditional source,
replace its weight $F_i$ by $w_i=F_i/(4\delta_i^8)$ to absorb
the guarded scalar square in the displayed multiplier.
Here $w_i$ is rational and depends only on the retained input,
square-abscissa data, and slope. In particular it has the same value
at every point-algebra character. All divisions and the displayed
nonzero square multipliers are guarded in the point stage.

Call a rational input \emph{prepared} if it has square reduced
denominator, lies in $(0,17)$, is a square in each of
$\Q_2,\Q_3,\Q_{17}$, and satisfies
\[
v_2(b)\ge3,\qquad v_3(b)\ge1,\qquad v_{17}(b)\ge2.
\]

\begin{corollary}\label{cor:nativeunits}
At an actual guarded point-stage tuple with prepared detector inputs,
each $17(17-b_i)$ is a unit of the whole point algebra.
\end{corollary}
\begin{proof}
At the selected rational character the square multiplier
in~\eqref{eq:detector} is a nonzero rational number. Preparation
therefore implies that $v_{17}(w_i)$ is even. Apply
Lemma~\ref{lem:units} separately to both inputs. The square multipliers
are units because the native point-stage guards hold at every
character.
\end{proof}

\begin{lemma}\label{lem:valuation}
Let $b\in\Q^\times$. If $17(17-b)$ is a rational square, then
$v_\ell(b)$ is not a negative odd integer at any prime $\ell\ne17$.
\end{lemma}
\begin{proof}
If $v_\ell(b)<0$ and $\ell\ne17$, the valuations of $17$ and $b$
are distinct, so
$v_\ell(17(17-b))=v_\ell(b)$. The valuation of a nonzero square
is even. The zero-complement case gives $b=17$ directly.
\end{proof}

Fix a coefficient member of the general source, or of either
conditional source, in the sufficient soundness range of its
corresponding theorem. Retain the source's remaining arithmetic
premises. The following is the additional arithmetic assertion relevant
to this reduction.

\begin{definition}[Native rational-complement realization]
\label{def:target}
For every integer input covered by the source completeness statement,
there is an actual rational point-stage tuple, with its native root
and auxiliary elliptic indices, for which all point-stage guard norms
are nonzero, both detector inputs are prepared, and
\[
                    17(17-b_0),\quad17(17-b_1)\in\Q^{\times2}
\]
at the selected rational character. For the conditional sources the
completeness range consists of the nonzero integers; the final input
factor includes zero.
\end{definition}

This assertion concerns the unmodified detector functions
\eqref{eq:detector}, evaluated at actual source tuples. It is a global
rational condition. Neither local squareness at finitely many places
nor the existence of freely chosen coefficient inputs implies it.

\begin{theorem}\label{thm:reduction}
If Definition~\ref{def:target} holds for the chosen source, its pole
stage can be replaced by two finite square choices using no retained
pole coordinates. The guarded one-coordinate finite-algebra encoding
then gives an integer polynomial defining $\Z$ with the following
total numbers of variables:
\begin{center}
\begin{tabular}{lrrr}\toprule
Source & Point pairs & Retained coordinates & Total\\\midrule
General &9&$1+18+1=20$&21\\
Either conditional &3&$1+6+1=8$&9\\\bottomrule
\end{tabular}
\end{center}
The input is included. The native realization assertion is not a
consequence of the guard lemmas.
\end{theorem}
\begin{proof}
Retain the input, all point pairs, and the slope. Let $R$ be the
integer polynomial ring in these coordinates, and let $\mathcal B_R$
be the free rank-$4^p$ point algebra. Clearing the native guarded
denominators gives $b_i=N_i/D_i$ with $N_i,D_i\in\mathcal B_R$.
The denominators are included among the guards. Adjoin two finite
generators by
\begin{equation}\label{eq:newalgebra}
\mathcal A_R=
\mathcal B_R[q_0,q_1]/
\bigl(q_i^2-17D_i(17D_i-N_i):i=0,1\bigr).
\end{equation}
This algebra is free of rank $4^{p+1}$. At a rational character with
$D_i\ne0$, the rational number $q_i/D_i$ is a square root of
$17(17-b_i)$, and every such root gives a character extending the
point character.

Include the norms of $17D_i-N_i$ in the scalar guard, together with
all retained point-stage guards and denominator guards. At a rational
base tuple where this scalar guard is nonzero, the point algebra is
etale and both new quadratic radicands are units. Thus all hypotheses
of the separator and determinant theorem in~\cite[Section 6]{Unit}
hold. Its polynomial envelopes and multiplication matrices use only
the retained coordinates. They introduce no further rational
parameters. With a single encoder coordinate $Y$, the resulting
integer determinant polynomial has a rational zero exactly when
\eqref{eq:newalgebra} has a rational character satisfying the guards.
If the scalar guard vanishes, that determinant is identically one.

For soundness, its rational character supplies the same native point
data and tests, together with both rational square complements.
Lemma~\ref{lem:valuation} gives the no-negative-odd-valuation
implication used by the detector soundness proofs of~\cite{General,Conditional}, outside their
exceptional set containing $2,3,17$. The remaining local extraction
and height arguments are unchanged. Thus the rational input is an
integer. For a conditional source, multiply the determinant by $X$
to include zero exactly as in its original source.

For completeness, Definition~\ref{def:target} supplies the actual
point tuple and both square roots. Corollary~\ref{cor:nativeunits}
makes the new guard factors nonzero at every point-algebra character.
The original point-stage guards are already nonzero by hypothesis.
Both radicands in~\eqref{eq:newalgebra} are therefore units, and the
separator theorem supplies the rational encoder coordinate. There are
$2p+2$ retained coordinates and one encoder, giving $2p+3$ in total.
\end{proof}

The two new finite roots replace the two former finite pole ordinates.
The four former retained coordinates $(\tau_0,s_0,\tau_1,s_1)$ have
all been removed. The proof requires neither an additional common
parameter nor any pole-guard density assertion. It retains every
point-stage guard. No degree or coefficient bound, and no established
variable record, is asserted without the realization hypothesis.

\section{Prepared free inputs at a common pole parameter}
The next result describes coefficient-level points on the rational
surface of~\cite{Quadric}.

\begin{proposition}\label{prop:preparation}
Let $w>0$ be rational and a square in $\Q_2$, $\Q_3$, and
$\Q_{17}$. Put $D=68w$ and fix $\tau\in\Q^\times$ with
$\tau^2\ne D$. There are infinitely many rational $h>0$, tending
to zero, for which
\[
b=\frac{68Dh^2}{(1+Dh^2)^2},\qquad
s=17\frac{1-Dh^2}{1+Dh^2},\qquad
V=\frac{1+Dh^2}{34h}
\]
satisfy $b=4w/V^2$, $17(17-b)=s^2$, and $b$ is prepared.
With
\[
e=\frac{\tau^2-D}{\tau^2+3D},\qquad
p=-\frac{4D\tau}{\tau^2+3D},\qquad q=ps^2,
\]
both pole equations
\[
p^2=D(1-e)(1+3e),\qquad
q^2=D(s^2-eC)(s^2+3eC),\quad C=17(17-b),
\]
hold with $pqs\ne0$. A finite family of such weights admits these
choices simultaneously at any one common admissible $\tau$.
\end{proposition}
\begin{proof}
The identities follow by substitution. Let $S$ contain $2,3,17$ and
the primes where $v_\ell(w)$ is odd. Choose integers $N_\ell\ge0$
such that
\[
v_\ell(D)+2N_\ell>0,\qquad
v_\ell(68D)+2N_\ell\ge r_\ell,
\]
where $(r_2,r_3,r_{17})=(3,1,2)$ and the other $r_\ell$ are zero.
Set $H=\prod_{\ell\in S}\ell^{N_\ell}$ and take $h=H/M$ with
positive $M\equiv1\pmod{\prod_{\ell\in S}\ell}$ tending to
infinity. At each $\ell\in S$, $1+Dh^2$ is a unit and the required
nonnegative valuation bound holds. Also
\[
\frac bw=\left(\frac{68h}{1+68wh^2}\right)^2,
\]
which gives the three local square conditions. Outside $S$, all
valuations of $b$ are even. Hence every negative valuation is even,
and the reduced denominator is a square. For sufficiently large $M$,
$0<b<17$ and $0<Dh^2<1$. The displayed $s,p,q$ are nonzero.
The function $b(h)$ is nonconstant, so these choices give infinitely
many distinct $b$. Apply the construction independently to each
weight for the finite-family assertion.
\end{proof}

Proposition~\ref{prop:preparation} does not impose
$V=\Tr(\mathcal W^{-1})$. Establishing that equality for actual
native tuples, simultaneously for both detectors and with their
point-stage guards, remains the arithmetic requirement in
Definition~\ref{def:target}.

\begin{thebibliography}{9}
\bibitem{Unit}
\emph{A single norm determinant for the 25- and 13-variable constructions},
project source, 14 November 2026, version 28,
\nolinkurl{papers/finite_algebra/unit_norm_20261114_v28.tex}.
\bibitem{General}
\emph{A defining polynomial in twenty-five variables: combined scalar
and comparison conditions}, project source, 10 November 2026, version 19,
\nolinkurl{papers/few_variables25/few_variables25_merged_20261110_v19.tex}.
\bibitem{Conditional}
\emph{Thirteen Variables from Direct Elliptic Index Comparison}, project source,
10 November 2026, version 22,
\nolinkurl{papers/conditional_few_variables/direct_input13_20261110_v22.tex}.
\bibitem{Quadric}
\emph{A rational quadric model for a varying-input pole surface}, project source,
16 November 2026, version 44,
\nolinkurl{papers/pole_parameters/varying_input_quadric_20261116_v44.tex}.
\end{thebibliography}
\end{document}
